\chapter[Introducción]{Introducción}
\label{cp:introduccion}
% Extensión estimada: sin límite; la ficha orienta 6--7 páginas (ver fichas/01-introduccion.md y guía §5).
% Fuentes: Plan de trabajo (Introducción y objetivos), Propuesta (secciones "Introducción y Propósito" y
% "Marco de la Problemática"), PLAN_TFG_GAVIOLA_FACUNDO.pdf y la revisión del tutor del 28/09/2026
% (tareas_revision_cap1-2.md, I1--I8). Sin hipótesis (D1): solo preguntas de investigación y objetivos.

Este capítulo presenta el problema que motiva el trabajo, las preguntas de investigación y los objetivos, y
describe cómo se organiza el documento.

\section{Contexto y motivación}
\label{sec:intro_contexto}

El cáncer de piel es uno de los tipos de cáncer más frecuentes del mundo. La última estimación de GLOBOCAN,
publicada en 2026, atribuye al año 2024 1\,138\,314 casos nuevos de cáncer de piel no melanoma y 337\,969 de
melanoma~\cite{sung2026global}. El melanoma es menos de una cuarta parte de los diagnósticos, pero causa casi
tantas muertes como el resto de los cánceres de piel juntos, 58\,398 frente a 58\,721.

El desenlace del melanoma depende del momento en que se lo diagnostica. La supervivencia relativa a cinco años
es de 100.0\,\% cuando se detecta localizado, baja a 76.0\,\% cuando alcanzó los ganglios regionales y a
34.0\,\% cuando hay metástasis a distancia~\cite{nci2026melanoma}. Por eso importan las herramientas de apoyo al
diagnóstico que permitan detectarlo a tiempo.

La detección temprana enfrenta varias limitaciones. Los especialistas son escasos: un relevamiento en 158
países estimó una media de 2.66 dermatólogos por cada 100\,000 habitantes, con un rango de 0.37 en los países de
ingreso bajo a 5.05 en los de ingreso alto~\cite{freeman2026global}. Donde faltan, la primera valoración queda en
manos de clínicos que necesitan la dermatoscopía, cuyo beneficio depende de una formación que no todos
tienen~\cite{dinnes2018dermoscopy}. Aun con esa formación, la lectura de una misma imagen varía entre
observadores según su nivel de experiencia~\cite{tschandl2019comparison}. A esto se suma la documentación, que
consume una parte grande del tiempo médico~\cite{arndt2017tethered}.

Esas limitaciones retrasan la atención en dos puntos, la interpretación de la lesión y la redacción de su
reporte, y de ahí surgen dos tareas computacionales: asignar una categoría a la lesión y redactar su reporte.

Para la primera se aplicaron métodos de aprendizaje profundo, como las \glspl{cnn} y los \glspl{vit} entrenados
de forma supervisada sobre imágenes etiquetadas~\cite{dosovitskiy2021image,liu2021swin}. Ya en 2017 una \gls{cnn}
alcanzó un desempeño comparable al de dermatólogos certificados~\cite{esteva2017dermatologistlevel}.
Sin embargo, las revisiones sistemáticas registran exactitudes altas en condiciones
controladas~\cite{magalhaes2024systematic,kreouzi2024learning}, y queda abierto si ese desempeño se conserva en
atención primaria~\cite{jones2022intelligence}. Además, un clasificador produce una categoría y no el texto que
la describe, de modo que la documentación queda sin asistencia.

Los \glspl{vlm} responden a esa limitación porque integran información visual y textual. Se entrenan con pares
imagen--descripción que ubican ambas modalidades en un espacio común, lo que permite clasificar una imagen sin
fijar de antemano las categorías~\cite{radford2021learning}. Los modelos generativos agregan un modelo de
lenguaje que produce texto libre condicionado a la imagen~\cite{liu2023visual}, y su adaptación al dominio
biomédico mostró que el procedimiento se traslada a imágenes médicas~\cite{li2023llavamed}. Pueden ser generales
o estar especializados en un dominio.

Quedan abiertos dos temas. El primero es la generalización. El desempeño de los sistemas de \gls{ia}
dermatológica cae en los tonos de piel menos representados en los datos de entrenamiento, que dominan los
conjuntos públicos~\cite{daneshjou2022disparities,groh2021neural}, y la población argentina tampoco se puede
suponer homogénea, porque el Registro Argentino de Melanoma Cutáneo documentó diferencias regionales en las
características de la enfermedad~\cite{loria2020cutaneous}. El segundo son los reportes diagnósticamente fieles,
que describen la lesión que la imagen muestra y afirman el diagnóstico correcto sin sostener lo que la imagen no
respalda. En este trabajo se denomina fidelidad diagnóstica a la coincidencia entre el diagnóstico que el
reporte afirma y el de referencia de la imagen, y calidad descriptiva a la adecuación del texto a la lesión que
la imagen muestra.

Las dos tareas están conectadas, porque el diagnóstico de la clasificación es lo que el texto debe afirmar. Un
texto fluido puede sostener un diagnóstico equivocado, y un diagnóstico acertado puede acompañarse de una
descripción que no corresponde a la lesión, por lo que no basta resolverlas por separado. El problema de
investigación es entonces el estudio conjunto de la clasificación de lesiones y de la generación de reportes con
\glspl{vlm}, con recursos de hardware limitados.

\section{Preguntas de investigación}
\label{sec:intro_preguntas}
El problema se traduce en dos preguntas, una por tarea, porque sus salidas, una categoría y un texto, se evalúan
con métricas que no se traducen entre sí.

\begin{enumerate}[label=\textbf{PI\arabic*.}]
  \item ¿Cómo se compara el desempeño de \glspl{vlm} generales y especializados con el de clasificadores
  visuales supervisados en la clasificación de lesiones de piel y en la detección de melanoma?
  \item ¿Qué fidelidad diagnóstica y qué calidad descriptiva alcanzan, por separado, los reportes que genera un
  \gls{vlm} ajustado a partir de la imagen?
\end{enumerate}

\section{Objetivos}
\label{sec:intro_objetivos}
El objetivo general es diseñar, implementar y evaluar un enfoque de aprendizaje profundo y de \glspl{vlm},
adaptados a partir de modelos preentrenados, para clasificar lesiones cutáneas y generar reportes
dermatológicos a partir de imágenes, y analizar su desempeño diagnóstico, la fidelidad diagnóstica de los
reportes y su viabilidad computacional con recursos de hardware limitados.

Los objetivos específicos son:
\begin{enumerate}[label=\textbf{OE\arabic*.}]
  \item Diseñar, implementar y evaluar, adaptando modelos preentrenados, clasificadores de lesiones cutáneas
  basados en \glspl{cnn}, \glspl{vit} y \glspl{vlm} para identificar posibles melanomas, y comparar los
  \glspl{vlm} generales y especializados con los clasificadores supervisados.
  \item Diseñar, implementar y evaluar un \gls{vlm}, adaptado de un modelo preentrenado, que integre información
  visual y textual para producir un diagnóstico preliminar y el reporte de la lesión, con evaluación separada de
  la fidelidad diagnóstica y la calidad descriptiva.
  \item Evaluar los modelos desarrollados en distintas condiciones: la generalización a otra población, la
  calibración de la confianza y el tratamiento de la incertidumbre, y el costo computacional.
\end{enumerate}

La Tabla~\ref{tab:intro_objetivos_preguntas} relaciona cada objetivo específico con la pregunta que contribuye a
responder y con los capítulos donde se lo describe y se lo evalúa.

\begin{table}[htbp]
  \centering
  \caption{Relación entre los objetivos específicos, las preguntas de investigación y los capítulos.}
  \label{tab:intro_objetivos_preguntas}
  \small
  \begin{tabularx}{\textwidth}{@{}p{0.29\textwidth}p{0.10\textwidth}XX@{}}
    \toprule
    Objetivo específico & Pregunta & Se diseña y describe en & Se evalúa y responde en \\
    \midrule
    OE1. Clasificadores de lesiones y de melanoma (\gls{cnn}, \gls{vit}, \gls{vlm}) & PI1 &
      Capítulo~\ref{cp:metodos}, Materiales y métodos &
      Capítulo~\ref{cp:res_clasif}, Resultados de clasificación \\
    \filasep
    OE2. \gls{vlm} para generar reportes & PI2 &
      Capítulo~\ref{cp:metodos}, Materiales y métodos &
      Capítulo~\ref{cp:res_gen}, Resultados de generación \\
    \filasep
    OE3. Evaluación en distintas condiciones & PI1 y PI2 &
      Capítulo~\ref{cp:metodos} (protocolo de evaluación) &
      Capítulos~\ref{cp:res_clasif} y~\ref{cp:res_gen} (resultados) y Capítulo~\ref{cp:discusion} (discusión y
      limitaciones) \\
    \bottomrule
  \end{tabularx}
\end{table}

\section{Organización del documento}
\label{sec:intro_organizacion}
Después de esta introducción, el documento se organiza en seis capítulos:

\begin{description}
  \item[Capítulo~\ref{cp:marco}, Marco teórico.] Este capítulo presenta los conceptos necesarios para leer el
  resto y el estado del arte.
  \item[Capítulo~\ref{cp:metodos}, Materiales y métodos.] Este capítulo presenta los conjuntos de datos, los
  modelos, el \emph{pipeline} y el protocolo de evaluación.
  \item[Capítulo~\ref{cp:res_clasif}, Resultados de clasificación.] Este capítulo presenta la comparación de las
  estrategias de clasificación y la evaluación fuera del conjunto de entrenamiento.
  \item[Capítulo~\ref{cp:res_gen}, Resultados de generación.] Este capítulo presenta la evaluación de los
  reportes generados y el costo de cómputo de las configuraciones del \emph{pipeline}.
  \item[Capítulo~\ref{cp:discusion}, Discusión.] Este capítulo presenta la interpretación conjunta de los
  resultados, la respuesta a las preguntas de investigación y las limitaciones.
  \item[Capítulo~\ref{cp:conclusiones}, Conclusiones.] Este capítulo presenta los hallazgos principales, el
  balance de los objetivos y el trabajo futuro.
\end{description}

Los apéndices reúnen el material propio que no entra en el cuerpo, y el anexo, las fichas de los conjuntos de
datos y de los modelos externos que el trabajo utiliza.
